\documentclass{article}
\usepackage[T1]{fontenc}
\usepackage{lmodern}
\usepackage{graphicx}
\usepackage{amsmath,amssymb}
\usepackage[a4paper,margin=2cm]{geometry}
\usepackage[absolute,overlay]{textpos}
\setlength{\TPHorizModule}{1mm}
\setlength{\TPVertModule}{1mm}
\usepackage[indonesian]{babel}
\usepackage{hyperref}
\setlength{\emergencystretch}{2em}
% unit: Halaman 11

\begin{document}

% === Halaman 11 (llm) ===

\section*{Contoh pembangkitan kunci oleh Alice}

\begin{itemize}
    \item Misalkan Alice ingin menghitung pasangan kunci publik dan kunci privatnya.
    \item Misalkan Alice memilih $p = 47$ dan $q = 71$ (keduanya prima), maka dapat dihitung:
    \[ n = p \times q = 3337 \]
    \[ \phi(n) = (p - 1) \times (q - 1) = (47 - 1)(71 - 1) = 46 \times 60 = 3220. \]
    \item Alice memilih kunci publik $e = 79$ (yang relatif prima dengan 3220 karena pembagi bersama terbesarnya adalah 1).
    \item Kunci publik Alice adalah $(e, n) = (79, 3337)$. Nilai $e$ dan $n$ dapat dipublikasikan kepada umum.
\end{itemize}

\end{document}
